% introduction.tex
% Introduction: motivation, technical gap, contribution, roadmap.

Portfolio construction, risk attribution, and stress testing often
take long synthetic time series of asset growth rates as inputs. A
good generator has to satisfy two things the downstream consumer
cares about. First, each single asset's synthetic series should look
like its real counterpart when you histogram it day by day: the tails
should be as heavy as real tails, the kurtosis should match, regimes
and jumps should show up with roughly the right frequency; in
statistical language this is the \emph{marginal distribution} of the
asset. Second, a stack of synthetic series across many assets should
co-move the way real markets co-move: on a day when the market
plunges, stocks with high factor loadings should plunge with it, tail
co-movement should look realistic, and factor structure should be
intact; this is the \emph{cross-sectional dependence} across assets.
The two requirements are largely orthogonal, and off-the-shelf
generators tend to handle only one of them. Per-asset generators
built from hidden Markov models~\cite{hamilton1989}, generalized
autoregressive conditional heteroskedasticity (GARCH)
families~\cite{bollerslev1986}, or empirical
bootstraps~\cite{politisRomano1994} handle the first: fit one of
these models to AAPL's growth-rate history and its synthetic paths
carry AAPL's tail thickness, kurtosis, and regime structure. But the
fit knows nothing about SPY, MSFT, or any other asset. Stack the
per-asset draws into a joint object, a matrix of synthetic
growth-rate time series $\{\teps_i\}_{i=1}^N$ with one $T$-length
path per asset, each $\teps_i$ matching its own asset's full-return
marginal, and the paths drift independently: on a day when synthetic
SPY plunges, synthetic AAPL has no reason to plunge with it. Real
AAPL would. Closing this gap without breaking the per-asset marginals
is the problem the rest of this paper addresses. The standard recipe for adding
factor structure is the single-index model (SIM) of
Sharpe~\cite{sharpe1963}, which writes each asset's growth rate as
$g_i = \alpha_i + \beta_i g_m + \eps_i$ with $\eps_i$ idiosyncratic.
Applied naively on top of a per-asset generator, the recipe sets
$\eps_i = \teps_i$ and inflates the marginal variance of $g_i$ by
exactly $\beta_i^2 \sigma_m^2$. For high-$\beta$ assets the inflation
is large enough to be visible in the empirical marginal: the body
widens, the kurtosis drops, and the heavy tails that the generator was
built to deliver are pushed out into a Gaussian-flavored bulge driven
by the market term. This defeats the purpose of fitting a heavy-tailed
generator in the first place. The dual failure mode,
$\eps_i \sim \mathcal{N}(0, \sigma_{\eps,\real}^2)$ with the residual
variance taken from the real ordinary least squares (OLS) fit,
recovers the SIM coefficients
exactly but erases the marginal structure entirely.

This paper describes a closed-form variance correction that sits
between the two failure modes. Setting $\eps_i = s_i\,\teps_i$ for a
scalar $s_i \in (0,1]$ chosen so that the variance algebra works out to
preserve the generator marginal, the scalar admits a closed form,
$s_i^2 = 1 - \beta_i^2 \sigma_m^2 / \sigma_{\gen,i}^2$
(Eq.~\ref{eq:scale}), defined whenever the bracketed quantity is
positive. When it is not (high $\beta_i$ paired with a low-volatility
generator, or a synthetic market more volatile than the calibration
market), $\beta_i$ is clipped downward against a floor that reserves a
configurable share of the marginal variance for idiosyncratic noise,
with a sign-preserving clip (Eq.~\ref{eq:clip}) that keeps defensive
assets defensive. A separate $R^2$-preserving branch
(Eq.~\ref{eq:r2-target}) handles index-tracking assets, where the
calibrated $R^2$ rather than the generator marginal is the quantity
worth preserving; SPY against itself recovers as the degenerate
$R^2 = 1$ limit with no special-casing. Branch selection is driven
entirely by the calibrated $R^2$, and each asset carries a persisted
construction flag (\texttt{hybrid}, \texttt{hybrid-clipped},
\texttt{r2-preserve}) so downstream consumers can audit which path
each asset took. The procedure generalizes to any new ticker without
external metadata.

In this study, we tested the construction on the $424$-ticker US
equity universe and SPY market path used by the companion
\texttt{JumpHMM.jl} validation paper of Alswaidan and
Varner~\cite{alswaidanVarner2026jdiq}, comparing three composition
recipes on the same generators and the same market path: the naive
composition, a Gaussian-residual SIM that recovers the calibrated
$(\alpha, \beta, R^2)$ but loses the marginal structure, and the
hybrid construction of Section~\ref{sec:method}. We measured SIM
recovery, marginal fidelity, and cross-sectional dependence, reporting
each per construction-flag branch. The hybrid construction held the
SIM-recovery error of the naive composition while holding the
marginal-fidelity error of the Gaussian-residual baseline across the
full range of empirical $\beta_i$; on the aggregate scorecard it
recovered $\beta$ to a median absolute error of $0.018$ and passed a
per-ticker Kolmogorov--Smirnov test against the real marginal at rate
$50.4\%$, compared to $4.5\%$ for the naive composition and $0.6\%$
for the Gaussian baseline. The broader implication is that the
marginal-versus-factor tradeoff long assumed in SIM composition is not
fundamental: a single closed-form scalar correction closes the gap,
and the correction generalizes to richer factor structures
(multi-factor SIM, lagged-factor SIM) without changing the underlying
algebra. The remainder of the paper formalizes the preservation criteria and
derives the variance correction and the $R^2$-preserving branch
(Section~\ref{sec:method}), describes
reports the experimental protocol and the recovery and fidelity
metrics (Section~\ref{sec:results}), and discusses caveats and
extensions (Section~\ref{sec:discussion}).
